\documentclass[10pt,a4paper]{amsart}
\usepackage[margin=19mm]{geometry}
\usepackage{amsmath,amssymb,mathtools}
\usepackage[scr=rsfs,cal=euler]{mathalfa}
\usepackage{xfrac}
\usepackage[unicode,hidelinks,bookmarks=false]{hyperref}
\newcommand{\ie}{i.e.\ }
\newcommand{\cf}{cf.\ }
\newcommand{\resp}{respectively\ }
\newcommand{\Subsection}{\subsection}
\providecommand{\nobreakdash}{\nobreak\hbox{-}}
\setlength{\emergencystretch}{2em}
\multlinegap0pt
\usepackage[shortlabels]{enumitem}
\usepackage{mathtools}
\newtheorem{theorem}{Theorem}
\newtheorem{proposition}{Proposition}
\newtheorem{lemma}{Lemma}
\def \B{\mathbb B}
\def \R{\mathbb R}
\def \e{\varepsilon}
\newcommand \vol[2][3]{\left|#2\right|_{#1}}
\title{On convex bodies with constant non-central sections}
\author{J. Haddad, D. Ryabogin}
\date{Source: arXiv:2605.00299v1 (CC BY 4.0)}
\begin{document}
\maketitle

\noindent\emph{Source and licence.} On convex bodies with constant non-central sections. Version arXiv:2605.00299v1 (CC BY 4.0), distributed by arXiv under the Creative Commons Attribution 4.0 International licence, http://creativecommons.org/licenses/by/4.0/, as recorded in the arXiv licence metadata for this exact version and shown on the abstract page https://arxiv.org/abs/2605.00299v1. Full licence notice: \texttt{licenses/arxiv-2605.00299-CC-BY-4.0.txt}.\par\smallskip

\noindent\textit{Editorial scope.} Counted unit: the article's proposition prop\_uniqueness\_cases (source0.tex lines 324--334). The article's complete original proof of it is delivered with it (source0.tex lines 335--395, one \texttt{proof} environment of 61 lines). No mathematics is written, repaired or re-derived: the statement and the proof are the article's own lines, in the article's own notation, and no retained line is substituted or re-flowed.  Retained uncounted as the source-local context the proof needs: lines 74--80; lines 126--129; lines 162--165; lines 197--200; lines 204--231; lines 252--255; lines 265--268; lines 278--283.  One declared adaptation: exactly 1 ancillary strings inside the retained spans are omitted (image-inclusion commands and, where the source repeats a label, the repeated label definitions) and no other line differs from the source; the figure environments, with their placement, captions and labels, are kept, and the package records the omitted strings in fidelity receipts bound to the affected bodies.  No retained line contains a citation, so the wrapper carries no bibliography and no bibliography entry.  The retained lines contain the article's own diagram code, which the wrapper typesets with the standard tikz packages; no image file is delivered and the package declares no asset dependency.  Editorial formatting only, in addition: an AMS \texttt{amsart} preamble that restates the article's own declarations and macros used by the retained lines, namely the environments \texttt{theorem}, \texttt{proposition}, \texttt{lemma} on separate counters, together with the standard packages those lines need.  The retained environments print in this document as Theorem 1, Proposition 1, Lemma 1 in order of appearance, whereas the article prints the same items with the numbering of its own full document.  The complete native source of the article as distributed in the arXiv e-print for this exact version is bundled unchanged as \texttt{sources/199550/source0.tex}.
\begin{theorem}
	\label{thm_main}
	There exists an infinite set $\mathbb A \subseteq \R$ of positive Hausdorff dimension with the following property:

	Let $C \subseteq \R^4$ be a symmetric convex body of revolution containing the centered unit Euclidean ball $\B_4$ in its interior, such that for every hyperplane $H$ tangent to $\B_4$, $\vol[3]{C \cap H} = A$ where $A$ is a constant independent of $H$.
	If $A \in \mathbb A$ then $C$ is a Euclidean ball.
\end{theorem}

\begin{equation}
	\label{def_K}
	K = C \cap \langle e_1, e_2 \rangle. %(\R^2 \times \{(0,0)\}).
\end{equation}

\begin{equation}
	\label{eq_line}
	y = s x + h(s), \text{ where } h(s) = \sqrt{1+s^2}.
\end{equation}

	\begin{equation}
		\label{eq_def_Ps}
		P_{s}(x) = -\frac 23 \frac{\sqrt{1+s^2}}{s} x^3 - \frac{1+2s^2}{s^2} x^2 -2 \frac{\sqrt{1+s^2}}{s} x.
	\end{equation}

Consider $s \in \R$ and a tangent line $L_s$ with slope $s$ which touches the ball at a point $(t(s),d(s) )$ in the upper half-plane (see Figure \ref{fig_critical}), where $t(s), d(s)$ are functions of $s$.
For $s\neq 0$, the intersection point of $L_s$ with the $X$ axis is a point $(z(s),0)$, where $z(s)$ is also a function of $s$.
The functions $t(s),z(s)$ have very special properties, and will be used extensively in this section.
Observe that $z(s)$ and $t(s)$ have the same sign, which is opposite to the sign of $s$.
	\begin{figure}
		\caption{The values $t(s)$ and $z(s)$ are the critical points of the polynomial $P_s$, whose graph is depicted in blue.}
		\label{fig_critical}
		
	\end{figure}
Elementary computations show that
\begin{align}
	\label{eq_t}
	t(s) = -h'(s) = - \frac {s}{\sqrt{1+s^2}},\\
	\label{eq_z}
	z(s) = -h(s)/s = - \frac {\sqrt{1+s^2}}{s},
\end{align}
where $h$ is as in \eqref{eq_line}.
The following proposition is elementary.
\begin{proposition}
	\label{prop_zt}
	The functions $t(s), z(s)$ satisfy the following properties:
	\begin{enumerate}[label = (\alph*), ref = \alph* ]
		\item $t(s)$ and $z(s)$ are the only two critical points of $P_{s}$.
		\item $t(s)$ is always a local maximum of $P_{s}$ while $z(s)$ is always a local minimum.
		%\item $z(s),t(s)$ always have the same sign.
		\item $z(s) t(s) = 1$. \label{eq_zt_product}
	\end{enumerate}
\end{proposition}

\begin{equation}
	\label{eq_middle}
	x < t(x,y) < v.
\end{equation}

\begin{equation}
	\label{eq_billardeq}
	P_{s}(v) = - \frac{\frac 43 (R^2-1)^{3/2}}{s(1+s^2)} + P_{s}(x) ,
\end{equation}

\begin{lemma}
	\label{lem_side}
	Let $M \subseteq \R^2$ be a convex body, symmetric with respect to the $X$ and $Y$ axes, containing $\B_2$ in its interior, and let $X(M)$ be the largest $X$-coordinate of a point in $M$.
	Then for every $p \in \partial M$ with $\e(p)=1$ and $s(p)<0$, 
	the $Y$-coordinate of $T_M(p)$ has the same sign as $s(p) + s(X(M),0)$.
\end{lemma}

\begin{proposition}
	\label{prop_uniqueness_cases}
	Let $K$ be defined by \eqref{def_K} where $C$ is as in Theorem \ref{thm_main}.
	Let $p \in S_R \cap \partial K$ be such that one of the following three conditions is satisfied:
	\begin{enumerate}[label = (\alph*), ref = \alph* ]
		\item \label{cond_region} $s(p) \geq 0$
		\item \label{cond_apriori} $3+2s(p) \sqrt{R^2-1} < 0$,
		\item \label{cond_interval} $s(p)$ is not between $-s(X(K),0)$ and $-s(R,0)$.
	\end{enumerate}
	Then $T_R(p) \in \partial K$.
\end{proposition}

\begin{proof}
	Since $K = -K$, $s(-p) = s(p)$ and $T_K(-p) = - T_K(p)$, we may assume always that $\e(p) = 1$.
	Let $(x,y) = p$.
	Take $(v,w) = T_K(x,y)$ and $(v',w') = T_R(x,y)$.
	Our goal in each condition is to prove that $T_K(p) = T_R(p)$, which implies that $T_R(p) \in \partial K$.
	This can be achieved by showing that $v = v'$, since both $(v,w)$ and $(v',w')$ belong to the line $L_p$.

	We begin with condition \eqref{cond_region}.
	If $s(p) = 0$ then it is clear that $v = v' = -x$ by the symmetry of $K$ with respect to the $Y$ axis.
	Now assume $s(p) > 0$.
	The third-degree polynomial $P_{s(p)}$ has two critical points, $t(p), z(p)$, of which $t(p)$ is a local maximum, and $z(p)$ is a local minimum.
	We recall that $x < t(p) < v,v'$ (see \eqref{eq_middle}), and by Property \eqref{eq_zt_product} in Proposition \ref{prop_zt}, we have $z(p) < -1 < t(p) <  0$.
	But this implies that $P_{s(p)}$ has no critical points in the interval $(t(p), \infty)$ making it an injective function (see Figure \ref{fig_critical}, left).
	Since both $v,v'$ are solutions of \eqref{eq_billardeq} in that interval, we conclude that $v=v'$ which means $T_R(p) = T_K(p)$, so the first part of the proof is done.

	Now let us consider condition \eqref{cond_apriori}.
	From the first part of the proof, we may assume $s(p) < 0$, which implies $0 < t(p) < z(p)$. From \eqref{eq_def_Ps} we see that
	\[\lim_{v \to \infty} P_{s(p)} (v) = +\infty.\]

	We claim that condition \eqref{cond_apriori} implies 
	\begin{equation}
		\label{eq_uniqueness_condition}
		P_{s(p)} (x) - \frac{\frac 43 (R^2-1)^{3/2}}{s(p)(1+s(p)^2)} > P_{s(p)} (t(p)).
	\end{equation}
	Given this claim, since $t(p)$ is the unique local maximum of $P_{s(p)} $, the polynomial $P_{s(p)} $ cannot have more than one pre-image of a point larger than $P_{s(p)} (t(p))$ (see Figure \ref{fig_critical}, right).
	Then again equation \eqref{eq_billardeq} has only one solution $v$ which forces $T_K(p) = T_R(p)$.
	Now let us prove the claim:

	Put $s = s(p)$. Since $p = (x,y) \in S_R$, we have
	\[
		x^2 + (x s + \sqrt{1+s^2})^2 = R^2.
	\]
	Solving for $x$ and considering $x < t(s) = - \frac s {\sqrt{1+s^2}}$ we find the value of $x$,
	\[
		x = - \frac s {\sqrt{1+s^2}} - \sqrt{\frac{R^2-1}{s^2+1}}.
	\]
	Replace this value of $x$ in \eqref{eq_uniqueness_condition} to obtain
	\begin{align}
		\frac{2 s^2}{3 (s^2+1)}-\frac{2 s^2+1}{s^2+1}+2 <
		-\frac{4 (R^2-1)^{3/2}}{3 s (s^2+1)}
		+\frac{2 \sqrt{s^2+1} }{s} \left(\sqrt{\frac{R^2-1}{s^2+1}} +\frac{s}{\sqrt{s^2+1}}\right) \\
		+\frac{2 \sqrt{s^2+1}}{3 s} \left(\sqrt{\frac{R^2-1}{s^2+1}}+\frac{s}{\sqrt{s^2+1}}\right)^3
		+\frac{2 s^2+1}{s^2} \left(\sqrt{\frac{R^2-1}{s^2+1}}+\frac{s}{\sqrt{s^2+1}}\right)^2 \\
	\end{align}
	which after some computations simplifies to \eqref{cond_apriori}, as required.
	
	Finally let us consider condition \eqref{cond_interval}.
	First notice that $-s(X(K),0), -s(R,0) < 0$.
	As before, we may assume $s(p) < 0$ and $\e(p) = 1$, which imply $x < t(p) < z(p)$.
	Since $t(p), z(p)$ are the only two critical points of $P_{s(p)} $, there can be at most one solution $v$ of \eqref{eq_billardeq} in each of the intervals
	\begin{equation}
		\label{eq_intervals}
		(-\infty, t(p)), (t(p), z(p)), (z(p), \infty).
	\end{equation}

	Once again, the solution in $(-\infty, t(p))$ has to be discarded, by \eqref{eq_middle}.
	Also, since $z(p) s(p) + h(s(p)) = 0$, the sign of $w = v s(p) + h(s(p))$ equals that of $z(p) - v$.
	Condition \eqref{cond_apriori} implies that $s(p) + s(X(K),0)$ and $s(p) + s(R,0)$ have the same sign.
	By Lemma \ref{lem_side} applied to $M=K$ and to $M=R \B_2$ we deduce that the signs of $w$ and $w'$ coincide, meaning that either $v,v' < z(p)$ or $v,v' > z(p)$.
	Since the interval $(-\infty, t(p))$ was discarded, $v, v'$ belong to the same interval in \eqref{eq_intervals} and are thus equal.
\end{proof}

\end{document}
